% Copyright 2026 Sreeram Anil
% SPDX-License-Identifier: CC-BY-4.0
% =============================================================================
%  Notation and estimation preliminaries — foundational chapter
% =============================================================================
\chapter{Notation and estimation preliminaries}
\label{ch:prelim}

This chapter fixes the notation used throughout the report and assembles, in one place, the
estimation-theoretic material that the seventy estimator chapters of
Parts~II--XII draw on. Nothing here is new; the purpose is to make every later chapter
self-contained at the level of \emph{symbols} and to state once, carefully, the four ideas that
recur: the Bayesian filtering recursion and its Gaussian specialisation, the two ways of
approximating a nonlinear transformation (analytic linearisation and moment matching),
observability and its quantitative degradation, and the measures --- consistency, bounds,
robustness --- by which the benchmark of Part~XIII judges a filter. A reader who is fluent in
\cite{jazwinski1970,andersonmoore1979,barshalom2001,simon2006,sarkka2013} can skip to
\cref{sec:prelim-obs}, which contains the battery-specific observability material, and to
\cref{sec:prelim-symbols}, which is the symbol table.

\section{Typographic and notational conventions}
\label{sec:prelim-notation}

\paragraph{Vectors, matrices and scalars.} Vectors are bold lowercase, $\vx$, $\vu$, $\vy$;
matrices are bold uppercase, $\mF$, $\mH$, $\mP$; scalars are italic, $z$, $v$, $i$, $T$.
The $j$-th component of $\vx$ is the scalar $x_j$ and the $(i,j)$ entry of $\mA$ is $A_{ij}$.
$\mI$ (or $\mI_n$) is the identity, $\bm{0}$ the zero vector or matrix as the context requires,
$\bm{e}_j$ the $j$-th unit vector. $\mA^\T$ is the transpose, $\mA\inv$ the inverse,
$\tr(\mA)$ the trace, $\diag(a_1,\dots,a_n)$ the diagonal matrix with the listed entries and
$\diag(\mA)$ the vector of diagonal entries of $\mA$. For symmetric matrices $\mA\succ 0$ and
$\mA\succeq 0$ mean positive definite and positive semi-definite, and $\mA \succeq \mB$ means
$\mA-\mB\succeq 0$. $\real^{n}$ and $\real^{m\times n}$ are the real vector and matrix spaces,
$\norm{\vx}$ the Euclidean norm and $\norm{\vx}_{\mA}^2 = \vx^\T\mA\vx$ the weighted square
norm induced by $\mA\succ 0$.

\paragraph{Time.} All estimators in \estkit{} are discrete-time. The integer $k\in\{0,1,2,\dots\}$
is the sample index and $\dt$ the (constant) sample period, so that physical time is $t_k = k\dt$.
The benchmark uses $\dt = \SI{0.1}{\second}$ throughout, i.e.\ a \SI{10}{\hertz} measurement rate,
which is representative of a production battery-management system (BMS). A continuous-time
signal is written $x(t)$ and its sampled version $x_k = x(t_k)$. The shorthand
$\vy_{1:k} = \{\vy_1,\dots,\vy_k\}$ denotes a measurement history.

\paragraph{Estimates, priors and posteriors.} A hat denotes an estimate, $\xhat$. The
superscripts $^-$ and $^+$ distinguish the two quantities that exist at every sample: $\xminus_k$
is the \emph{prior} (predicted, \emph{a priori}) estimate, formed from $\vy_{1:k-1}$, and
$\xplus_k$ is the \emph{posterior} (updated, \emph{a posteriori}) estimate, formed from
$\vy_{1:k}$. The same convention applies to the error covariance, $\Pminus_k$ and $\Pplus_k$.
Estimation errors are $\tilde{\vx}^-_k = \vx_k - \xminus_k$ and
$\tilde{\vx}^+_k = \vx_k - \xplus_k$, so that
\begin{equation}
\Pminus_k = \E{\tilde{\vx}^-_k(\tilde{\vx}^-_k)^\T\mid\vy_{1:k-1}},\qquad
\Pplus_k = \E{\tilde{\vx}^+_k(\tilde{\vx}^+_k)^\T\mid\vy_{1:k}} .
\label{eq:prelim-Pdef}
\end{equation}
Where no confusion can arise the superscript is dropped and $\xhat_k$, $\mP_k$ mean the posterior.

\paragraph{Probability.} $\E{\cdot}$ is expectation, $\Var{\cdot}$ variance,
$\N{\bm{\mu}}{\bm\Sigma}$ the Gaussian distribution with mean $\bm\mu$ and covariance
$\bm\Sigma$, and $p(\cdot)$ a probability density. $\vw_k$ always denotes process noise and
$\vv_k$ measurement noise; $\mQ$ and $\mR$ are their covariances. The symbol
$\veps_k = \vy_k - \yhat_k$ is reserved for the \emph{innovation} (measurement residual) and
$\mS_k$ for its covariance.

\paragraph{Battery symbols.} $\soc\in[0,1]$ is the state of charge (SOC), $\ocv(\soc,T)$ the
open-circuit voltage, $v$ the terminal voltage, $i$ the cell current with the
\textbf{discharge-positive} sign convention used consistently in the library and this report,
$T$ the cell temperature in kelvin and $Q$ the cell capacity in ampere-hours. The reason for the
discharge-positive convention is historical (it is Plett's~\cite{plett2004ekf1,plett2015bms1})
and it makes the Coulomb-counting equation read $\soc_{k+1} = \soc_k - \eta i\dt/(3600Q)$;
every sign in \cref{ch:ecm,ch:spm} follows from it.

\paragraph{Code.} Identifiers, file names and option names taken verbatim from the library are
set in \code{typewriter} type; estimator names as registered in the benchmark are set as
\estname{EKF}, \estname{UKF}, and so on.

\section{The model class}
\label{sec:prelim-model}

Every estimator in \estkit{} is written against a single interface: a discrete-time,
input-driven, nonlinear model with additive noise,
\begin{align}
\vx_{k}   &= f(\vx_{k-1}, \vu_{k-1}) + \vw_{k-1}, & \vw_k &\sim \N{\bm 0}{\mQ_k},
  \label{eq:prelim-proc}\\
\vy_{k}   &= h(\vx_{k}, \vu_{k}) + \vv_{k},       & \vv_k &\sim \N{\bm 0}{\mR_k},
  \label{eq:prelim-meas}
\end{align}
where $\vx_k\in\real^{n_x}$ is the state, $\vu_k\in\real^{n_u}$ the input, $\vy_k\in\real^{n_y}$
the measurement, and $f:\real^{n_x}\times\real^{n_u}\to\real^{n_x}$,
$h:\real^{n_x}\times\real^{n_u}\to\real^{n_y}$ are continuously differentiable. The noises are
mutually independent white sequences, independent of the initial condition
$\vx_0\sim\N{\xhat_0}{\mP_0}$.

\begin{remark}[Inputs and scheduling variables]\label{rem:prelim-sched}
The vector $\vu_k$ carries two logically different kinds of quantity and the distinction matters
for the noise model. Some components are \emph{true inputs}: they excite the state and their
measurement error propagates into the state estimate. Others are \emph{scheduling variables}:
they do not excite the state but parameterise $f$ and $h$. For the battery cell of
\cref{ch:ecm}, $\vu = [\,i,\ T\,]^\T$: the current $i$ is a true input, while the temperature
$T$ merely schedules the Arrhenius-scaled resistances and time constants. Consequently the
process-noise covariance is built from the \emph{current}-sensor specification alone
(\cref{sec:ecm-estimator}), and a temperature-sensor error shows up as a slowly varying
parameter error, not as process noise. Writing $\vu$ for both is a deliberate simplification
of the software interface; the modelling consequence is stated here once so that later chapters
can refer to it.
\end{remark}

In software, a \emph{model} is any class that provides the dimensions $n_x, n_u, n_y$ and the
five functions $f$, $h$, $\mQ$, $\mR$ and (optionally) the Jacobians
\begin{equation}
\mF(\vx,\vu) = \frac{\partial f}{\partial \vx},\qquad
\mB(\vx,\vu) = \frac{\partial f}{\partial \vu},\qquad
\mH(\vx,\vu) = \frac{\partial h}{\partial \vx},
\label{eq:prelim-jac}
\end{equation}
together with a projection $\Pi$ onto a feasible set $\mathcal{X}\subset\real^{n_x}$
(\code{constrain}) and, for joint and dual estimation, a parameter vector
$\vtheta\in\real^{n_\theta}$ with accessors. If an analytic Jacobian is absent the library
substitutes a central-difference approximation with relative step
$\varepsilon_{\mathrm{mach}}^{1/3}$ \cite[\S 5]{golubvanloan2013}; \cref{ch:ecm} gives all three
Jacobians of the cell model in closed form, which is why the cell benchmark is cheap enough to
run seventy estimators over twelve scenarios.

\section{The Bayesian filtering recursion}
\label{sec:prelim-bayes}

\begin{definition}[Filtering, prediction, smoothing]\label{def:prelim-densities}
Given the model \eqref{eq:prelim-proc}--\eqref{eq:prelim-meas}, the \emph{filtering density} is
$p(\vx_k\mid\vy_{1:k})$, the \emph{one-step predictive density} is $p(\vx_k\mid\vy_{1:k-1})$, and
the \emph{smoothing densities} are $p(\vx_k\mid\vy_{1:N})$ for $N>k$.
\end{definition}

Because \eqref{eq:prelim-proc} is a Markov chain and \eqref{eq:prelim-meas} is conditionally
independent given the state, the filtering density obeys an exact two-step recursion
\cite{stratonovich1960,holee1964,jazwinski1970}. The prediction step is the
Chapman--Kolmogorov equation
\begin{equation}
p(\vx_k\mid\vy_{1:k-1}) = \int_{\real^{n_x}} p(\vx_k\mid\vx_{k-1})\,p(\vx_{k-1}\mid\vy_{1:k-1})\,
  \mathrm{d}\vx_{k-1},
\label{eq:prelim-ck}
\end{equation}
where $p(\vx_k\mid\vx_{k-1}) = \N{f(\vx_{k-1},\vu_{k-1})}{\mQ_{k-1}}$ evaluated at $\vx_k$ is the
transition kernel. The correction step is Bayes' rule,
\begin{equation}
p(\vx_k\mid\vy_{1:k}) = \frac{p(\vy_k\mid\vx_k)\,p(\vx_k\mid\vy_{1:k-1})}
 {\displaystyle\int p(\vy_k\mid\vx_k)\,p(\vx_k\mid\vy_{1:k-1})\,\mathrm{d}\vx_k},
\label{eq:prelim-bayes}
\end{equation}
with likelihood $p(\vy_k\mid\vx_k) = \N{h(\vx_k,\vu_k)}{\mR_k}$ evaluated at $\vy_k$. The
denominator of \eqref{eq:prelim-bayes} is the one-step \emph{evidence} $p(\vy_k\mid\vy_{1:k-1})$;
its product over $k$ is the marginal likelihood
\begin{equation}
p(\vy_{1:N}) = \prod_{k=1}^{N} p(\vy_k\mid\vy_{1:k-1}),
\label{eq:prelim-evidence}
\end{equation}
which is the quantity that multiple-model methods use to weight a bank of filters and that
maximum-likelihood tuning maximises over $(\mQ,\mR)$.

\eqref{eq:prelim-ck}--\eqref{eq:prelim-bayes} is a complete solution of the estimation problem
and is also completely useless in this form, because the integrals are intractable for all but
two model classes: the linear-Gaussian case, where the densities stay Gaussian and the recursion
closes on two moments (\cref{sec:prelim-kf}), and the finite-state case, where the integrals are
sums. Every chapter of this report is an answer to the question ``what do you do instead?'', and
the answers fall into four families:
\begin{enumerate}[nosep,leftmargin=*]
\item \textbf{Keep the Gaussian form, approximate the moments.} Linearise ($\Rightarrow$ EKF),
  match moments with a deterministic quadrature rule ($\Rightarrow$ UKF, CDKF, CKF, GHKF), or
  iterate the linearisation to a fixed point ($\Rightarrow$ IEKF, IPLF).
\item \textbf{Keep the recursion, approximate the density by samples.} Sequential importance
  resampling and its variants ($\Rightarrow$ particle, ensemble and Gaussian-sum filters).
\item \textbf{Abandon probability, design an error system.} Deterministic observers with
  provable error dynamics ($\Rightarrow$ Luenberger, sliding-mode, high-gain, interval).
\item \textbf{Solve an optimisation problem over a window.} Maximum \emph{a posteriori}
  estimation on a moving horizon ($\Rightarrow$ MHE, batch nonlinear least squares, smoothers).
\end{enumerate}

\section{The Gaussian case: the Kalman filter}
\label{sec:prelim-kf}

Specialise \eqref{eq:prelim-proc}--\eqref{eq:prelim-meas} to the linear time-varying model
\begin{align}
\vx_k &= \mF_{k-1}\vx_{k-1} + \mG_{k-1}\vu_{k-1} + \vw_{k-1}, \label{eq:prelim-lin-proc}\\
\vy_k &= \mH_k\vx_k + \vv_k. \label{eq:prelim-lin-meas}
\end{align}

\begin{theorem}[Kalman \cite{kalman1960}]\label{thm:prelim-kf}
For the model \eqref{eq:prelim-lin-proc}--\eqref{eq:prelim-lin-meas} with
$\vx_0\sim\N{\xhat_0}{\mP_0}$ and Gaussian, white, mutually independent $\vw_k,\vv_k$, the
filtering density is Gaussian for every $k$, $p(\vx_k\mid\vy_{1:k}) = \N{\xplus_k}{\Pplus_k}$,
and its moments obey
\begin{align}
\xminus_k &= \mF_{k-1}\xplus_{k-1} + \mG_{k-1}\vu_{k-1}, &
\Pminus_k &= \mF_{k-1}\Pplus_{k-1}\mF_{k-1}^\T + \mQ_{k-1},
  \label{eq:prelim-kf-tu}\\
\veps_k &= \vy_k - \mH_k\xminus_k, &
\mS_k &= \mH_k\Pminus_k\mH_k^\T + \mR_k, \label{eq:prelim-kf-innov}\\
\mK_k &= \Pminus_k\mH_k^\T\mS_k\inv, &&\notag\\
\xplus_k &= \xminus_k + \mK_k\veps_k, &
\Pplus_k &= (\mI - \mK_k\mH_k)\Pminus_k. \label{eq:prelim-kf-mu}
\end{align}
$\xplus_k$ is the conditional mean and therefore the minimum-mean-square-error (MMSE) estimate
among \emph{all} measurable functions of $\vy_{1:k}$.
\end{theorem}

\paragraph{Derivation by orthogonality.} The linear-MMSE argument is worth repeating because it
is the one that survives when the Gaussian assumption is dropped. Seek the estimator
$\xplus_k = \xminus_k + \mK_k\veps_k$ that minimises
$J(\mK_k) = \E{\norm{\vx_k-\xplus_k}^2} = \tr\Pplus_k$. Substituting
$\tilde{\vx}^+_k = \tilde{\vx}^-_k - \mK_k(\mH_k\tilde{\vx}^-_k + \vv_k)$ into
\eqref{eq:prelim-Pdef} and using $\E{\tilde{\vx}^-_k\vv_k^\T}=\bm 0$ gives the Joseph form
\begin{equation}
\Pplus_k = (\mI-\mK_k\mH_k)\Pminus_k(\mI-\mK_k\mH_k)^\T + \mK_k\mR_k\mK_k^\T,
\label{eq:prelim-joseph}
\end{equation}
valid for \emph{any} gain $\mK_k$. Differentiating $\tr\Pplus_k$ with respect to $\mK_k$ and
equating to zero,
$\partial \tr\Pplus_k/\partial\mK_k = -2\Pminus_k\mH_k^\T + 2\mK_k\mS_k = \bm0$, yields
$\mK_k = \Pminus_k\mH_k^\T\mS_k\inv$; substituting this optimal gain back into
\eqref{eq:prelim-joseph} collapses it to \eqref{eq:prelim-kf-mu}. Equivalently, the optimal gain
is the one that makes the posterior error orthogonal to the innovation,
$\E{\tilde{\vx}^+_k\veps_k^\T}=\bm 0$ --- the \emph{orthogonality principle}
\cite{kailath2000}. This derivation uses only second moments: the Kalman filter is the best
\emph{linear} estimator for any noise distribution with the given covariances, and the best
estimator outright only in the Gaussian case.

\begin{remark}[Why the Joseph form is used everywhere in \estkit{}]\label{rem:prelim-joseph}
\eqref{eq:prelim-kf-mu} is a difference of two positive semi-definite matrices and can lose
symmetry or definiteness through round-off; \eqref{eq:prelim-joseph} is a sum of two positive
semi-definite terms and cannot \cite{bucyjoseph1968,verhaegen1986}. More important for this
report, \eqref{eq:prelim-joseph} remains the \emph{correct} covariance when the gain has been
deliberately modified --- by a robust weight, a fading factor, an $H_\infty$ correction or a
constraint projection --- whereas \eqref{eq:prelim-kf-mu} does not. Every filter in
Parts~III--V that modifies $\mK_k$ therefore uses \eqref{eq:prelim-joseph}.
\end{remark}

\paragraph{Information form.} Define the information matrix $\bm{Y} = \mP\inv$ and the
information vector $\bm\xi = \mP\inv\xhat$. The measurement update of
\cref{thm:prelim-kf} becomes an addition,
\begin{equation}
\bm{Y}^+_k = \bm{Y}^-_k + \mH_k^\T\mR_k\inv\mH_k, \qquad
\bm{\xi}^+_k = \bm{\xi}^-_k + \mH_k^\T\mR_k\inv\vy_k,
\label{eq:prelim-info}
\end{equation}
by the matrix-inversion lemma
$(\mA+\mU\mC\bm{V})\inv = \mA\inv - \mA\inv\mU(\mC\inv+\bm{V}\mA\inv\mU)\inv\bm{V}\mA\inv$
\cite{kailath2000}. Three properties of \eqref{eq:prelim-info} are exploited later: it admits a
diffuse prior $\bm{Y}^-_0=\bm 0$ (total ignorance) which the covariance form cannot represent;
it costs $\mathcal{O}(n_x^2 n_y)$ rather than $\mathcal{O}(n_x^2 n_y + n_y^3)$ when
$n_y\gg n_x$; and the contributions of independent sensors simply add, which is the basis of
every distributed and federated architecture in Part~XI.

\paragraph{Steady state.} If the model is time-invariant and $(\mF,\mH)$ is detectable with
$(\mF,\mQ^{1/2})$ stabilisable, $\Pminus_k$ converges to the unique stabilising solution of the
discrete algebraic Riccati equation
\begin{equation}
\mP = \mF\left[\mP - \mP\mH^\T(\mH\mP\mH^\T+\mR)\inv\mH\mP\right]\mF^\T + \mQ,
\label{eq:prelim-dare}
\end{equation}
and the gain becomes constant \cite{andersonmoore1979}. The steady-state Kalman filter obtained
this way is, structurally, a Luenberger observer with a particular gain
(\cref{sec:prelim-observers}); \estkit{} solves \eqref{eq:prelim-dare} once at \code{init()} for
the \estname{SSKF} and for the Riccati-designed observer gains.

\paragraph{Smoothing.} Post-flight analysis does not have to be causal. The
Rauch--Tung--Striebel recursion \cite{rauch1965rts} obtains the smoothing densities of
\cref{def:prelim-densities} by a single backward sweep over the quantities the forward filter has
already stored,
\begin{align}
\bm{G}_k &= \Pplus_k\mF_k^\T\big(\Pminus_{k+1}\big)\inv, \label{eq:prelim-rtsgain}\\
\xhat^{s}_k &= \xplus_k + \bm{G}_k\left(\xhat^{s}_{k+1} - \xminus_{k+1}\right), \qquad
\mP^{s}_k = \Pplus_k + \bm{G}_k\left(\mP^{s}_{k+1} - \Pminus_{k+1}\right)\bm{G}_k^\T,
\label{eq:prelim-rts}
\end{align}
initialised at $\xhat^s_N = \xplus_N$, $\mP^s_N = \Pplus_N$. Because
$\mP^{s}_{k+1}\preceq\Pminus_{k+1}$, \eqref{eq:prelim-rts} can only reduce the covariance: a
smoother is never worse than the filter and is typically two to three times better in the middle
of a record. Part~X treats the extended, unscented and fixed-lag variants; the relevance here is
that the smoother defines the accuracy \emph{ceiling} against which the causal estimators of the
benchmark should be read.

\paragraph{Factored implementations.} \eqref{eq:prelim-kf-tu}--\eqref{eq:prelim-kf-mu} propagate
$\mP$, whose condition number is the \emph{square} of that of the underlying estimation problem.
On a fixed-point or single-precision target this matters, and the classical remedy is to
propagate a factor instead: a Cholesky factor $\mS$ with $\mP = \mS\mS^\T$, updated by QR
(square-root filters \cite{potter1963,kaminski1971}), or the $\mU\mD\mU^\T$ decomposition with
$\mU$ unit upper triangular and $\mD$ diagonal \cite{bierman1977,thornton1976ud}. Both are
algebraically identical to \cref{thm:prelim-kf} --- Part~V verifies that they reproduce the EKF
to $10^{-11}$ --- but they halve the dynamic range required of the arithmetic and make an
indefinite covariance structurally impossible. The number to compare against
$1/\varepsilon_{\mathrm{float}} \approx 8\times10^{6}$ when deciding whether the factored form is
necessary is the peak condition number of $\mP$ along the trajectory.

\section{Linearisation and the extended Kalman filter}
\label{sec:prelim-ekf}

The extended Kalman filter (EKF) applies \cref{thm:prelim-kf} to the first-order Taylor
expansions of $f$ and $h$ about the current estimate,
\begin{equation}
\mF_{k-1} = \left.\frac{\partial f}{\partial\vx}\right|_{\xplus_{k-1},\vu_{k-1}},\qquad
\mH_{k}   = \left.\frac{\partial h}{\partial\vx}\right|_{\xminus_{k},\vu_{k}} ,
\label{eq:prelim-ekf-jac}
\end{equation}
with the nonlinear functions retained in the mean propagation,
$\xminus_k = f(\xplus_{k-1},\vu_{k-1})$ and $\yhat_k = h(\xminus_k,\vu_k)$. Writing the second-order
expansion of $h$ about $\xminus_k$,
\begin{equation}
h(\vx) = h(\xminus_k) + \mH_k\tilde{\vx}^-_k
 + \tfrac12 \sum_{j} \bm{e}_j\,(\tilde{\vx}^-_k)^\T \bm{\nabla}^2 h_j\, \tilde{\vx}^-_k
 + \mathcal{O}(\norm{\tilde{\vx}^-_k}^3),
\label{eq:prelim-ekf-2nd}
\end{equation}
and taking expectations shows that the EKF makes two distinct errors. The predicted measurement
is biased by $\tfrac12\sum_j \bm e_j \tr(\bm\nabla^2 h_j\,\Pminus_k)$, and the innovation
covariance \eqref{eq:prelim-kf-innov} omits the curvature contribution
$\tfrac12\tr(\bm\nabla^2 h_j\Pminus_k\bm\nabla^2 h_l\Pminus_k)$. The first error is
$\mathcal{O}(\norm{\tilde{\vx}}^2)$, the second $\mathcal{O}(\norm{\tilde{\vx}}^3)$; both are
\emph{one-sided}, in that the omitted covariance term is positive semi-definite. The EKF is
therefore systematically \emph{over-confident}, and over-confidence is not a cosmetic defect: an
under-sized $\mS_k$ means an over-sized $\mK_k$ at the wrong time and an under-sized $\mK_k$
later, which is exactly the mechanism by which the EKF locks in an initial error on a curved
$\ocv(\soc)$ characteristic. The second-order filter restores the missing terms explicitly; the
sigma-point filters of the next section restore them implicitly.

\begin{assumption}[Local linearity]\label{as:prelim-lin}
Over a region of the state space whose diameter is set by $\mP$, $f$ and $h$ are well
approximated by \eqref{eq:prelim-ekf-jac}.
\end{assumption}

\cref{as:prelim-lin} is the single approximation behind every filter in Part~III that carries the
prefix ``extended'', and almost every failure recorded in Part~XIII can be traced to its
violation. It is a joint condition on the model \emph{and} on the uncertainty: a strongly curved
$h$ is harmless under a tight prior, and a mildly curved $h$ is fatal under a loose one.

\section{Moment matching and sigma-point rules}
\label{sec:prelim-sigma}

The Gaussian filters of Part~III that are not ``extended'' replace linearisation by numerical
integration. The observation is that \eqref{eq:prelim-ck}--\eqref{eq:prelim-bayes}, restricted to
Gaussian densities, needs only three integrals of the form
\begin{equation}
\E{g(\vx)} = \int_{\real^{n_x}} g(\vx)\,\N{\vx;\,\bm\mu}{\bm\Sigma}\,\mathrm{d}\vx ,
\label{eq:prelim-gaussint}
\end{equation}
namely $g = f$, $g = h$ and the corresponding outer products \cite{ito2000gaussian}. Any rule
that evaluates \eqref{eq:prelim-gaussint} accurately will do.

\begin{definition}[Sigma-point rule]\label{def:prelim-sigma}
A sigma-point (cubature) rule of order $m$ is a set of points
$\bm{\chi}^{(j)} = \bm\mu + \bm{L}\bm{\zeta}^{(j)}$, where $\bm{L}\bm{L}^\T = \bm\Sigma$, with
weights $w_m^{(j)}$ and $w_c^{(j)}$ such that
$\sum_j w_m^{(j)} g(\bm\chi^{(j)})$ reproduces \eqref{eq:prelim-gaussint} exactly for every
polynomial $g$ of total degree $\le m$.
\end{definition}

The unscented transform \cite{julier2000newmethod,julier2004unscented,wan2000ukf} uses
$2n_x+1$ points with $\bm\zeta^{(0)}=\bm0$ and
$\bm\zeta^{(j)}=\pm\sqrt{n_x+\lambda}\,\bm{e}_j$, $\lambda = \alpha^2(n_x+\kappa)-n_x$,
\begin{equation}
w_m^{(0)} = \frac{\lambda}{n_x+\lambda},\quad
w_c^{(0)} = w_m^{(0)} + (1-\alpha^2+\beta),\quad
w_m^{(j)} = w_c^{(j)} = \frac{1}{2(n_x+\lambda)},
\label{eq:prelim-utweights}
\end{equation}
and is exact to third order for the mean of any $g$ and captures the second-order term of the
covariance. The spherical-radial cubature rule \cite{arasaratnam2009ckf} drops the centre point,
uses $2n_x$ equally weighted points at radius $\sqrt{n_x}$, and is the $\alpha=1,\beta=0,\kappa=0$
limit of \eqref{eq:prelim-utweights}; the central-difference filter \cite{norgaard2000} obtains
the same accuracy from Stirling interpolation; Gauss--Hermite rules reach arbitrary order at a
cost of $m^{n_x}$ points. The essential point for this report is that all of them are
\emph{derivative-free}, so they never need $\mF$ or $\mH$, and all of them propagate the prior
covariance through the true nonlinearity, so the innovation covariance they produce is larger
than the EKF's --- which is why they survive the scenarios in which the EKF does not.

\begin{remark}[Statistical linearisation]\label{rem:prelim-statlin}
A sigma-point filter can equivalently be read as an EKF whose Jacobian has been replaced by the
statistical linear regression of $g$ over the prior, $\mH^{\mathrm{slr}} =
\bm{P}_{xy}^\T\bm\Sigma\inv$ with $\bm{P}_{xy}$ the cross-covariance estimated from the points.
This reading explains both the robustness (the regression averages the nonlinearity over the
prior rather than sampling its tangent at one point) and the failure mode (when the prior is
very wide, the regression slope can be far from the local slope at the truth). The iterated
posterior linearisation filter makes the regression point the posterior and iterates.
\end{remark}

\section{The optimisation view}
\label{sec:prelim-map}

The recursions above propagate a density. An equivalent and sometimes more convenient view is to
write down the joint posterior over a trajectory and maximise it. Taking logarithms of
\eqref{eq:prelim-ck}--\eqref{eq:prelim-bayes} over a window of $N$ samples and dropping constants,
the maximum-\emph{a-posteriori} (MAP) trajectory estimate solves
\begin{equation}
\min_{\vx_{0:N}}\ \
\underbrace{\norm{\vx_0 - \xhat_0}^2_{\mP_0\inv}}_{\text{arrival cost}}
+ \sum_{k=1}^{N}\norm{\vx_k - f(\vx_{k-1},\vu_{k-1})}^2_{\mQ\inv}
+ \sum_{k=1}^{N}\norm{\vy_k - h(\vx_k,\vu_k)}^2_{\mR\inv}
\quad \text{s.t. } \vx_k\in\mathcal{X}.
\label{eq:prelim-map}
\end{equation}
Three facts connect \eqref{eq:prelim-map} to the rest of the chapter. In the linear-Gaussian case
its solution coincides exactly with the Rauch--Tung--Striebel smoother
\eqref{eq:prelim-rtsgain}--\eqref{eq:prelim-rts}, and its last element with the Kalman filter, so
nothing is lost by the change of view. In the nonlinear case it is \emph{not} the same as the
EKF: \eqref{eq:prelim-map} relinearises at the optimum rather than at the prediction, which is
precisely the difference between the EKF and the iterated EKF, and between filtering and
full-information estimation. And it is the only formulation in which the constraint set
$\mathcal{X}$ enters as a constraint rather than as an after-the-fact projection --- the reason
moving-horizon estimation (Part~IX) handles hard bounds on $\soc$ natively while the Kalman family
must clip. The price is the cost of the optimisation: truncating the window to a fixed horizon
and summarising the discarded past by an arrival cost is what makes it tractable
\cite{rao2003mhe,rawlings2017mpc}.

\cref{tab:prelim-roadmap} summarises how the four families answer the same question.

\begin{table}[H]\centering
\caption{The approximation families and what each does to the exact recursion
\eqref{eq:prelim-ck}--\eqref{eq:prelim-bayes}.}
\label{tab:prelim-roadmap}
\small
\begin{tabular}{@{}p{30mm}p{52mm}p{50mm}@{}}\toprule
Family & Approximation & Fails when \\ \midrule
Analytic linearisation & first-order Taylor \eqref{eq:prelim-ekf-jac}; density stays Gaussian &
 prior wide relative to the curvature of $f,h$ (\cref{as:prelim-lin}) \\
Moment matching & deterministic quadrature \eqref{eq:prelim-gaussint}; density stays Gaussian &
 posterior genuinely multi-modal or heavy-tailed \\
Monte Carlo & density represented by weighted samples &
 $n_x$ large, or likelihood far sharper than the proposal \\
Deterministic observer & no density; error dynamics shaped directly &
 an uncertainty figure is required, or the disturbance is unbounded \\
Optimisation & MAP over a window \eqref{eq:prelim-map} &
 the horizon must be short enough to compute, losing older information \\
\bottomrule\end{tabular}
\end{table}

\section{Observability}
\label{sec:prelim-obs}

Observability is the property that decides whether an estimation problem is solvable at all,
and its quantitative degradation is the single best predictor of benchmark performance in
Part~XIII.

\subsection{Linear systems}
For the time-invariant pair $(\mF,\mH)$ the $n$-step observability matrix and observability
Gramian over a horizon of $N$ samples are
\begin{equation}
\mathcal{O}_N = \begin{bmatrix}\mH\\ \mH\mF\\ \vdots\\ \mH\mF^{N-1}\end{bmatrix}
\in\real^{N n_y\times n_x},
\qquad
\mathcal{W}_N = \mathcal{O}_N^\T\mathcal{O}_N = \sum_{j=0}^{N-1}(\mF^{j})^\T\mH^\T\mH\mF^{j}.
\label{eq:prelim-obsmat}
\end{equation}
The pair is observable if and only if $\operatorname{rank}\mathcal{O}_{n_x}=n_x$, equivalently
$\mathcal{W}_{n_x}\succ 0$ \cite{kalman1960controllability,chenct1999linear,kailath2000}. Rank,
however, is a yes/no statement about exact arithmetic and tells us nothing about a real filter.
The useful object is the spectrum of $\mathcal{W}_N$: the energy that an initial state
$\vx_0$ contributes to the noise-free output over the horizon is
$\sum_{j}\norm{\vy_j}^2 = \vx_0^\T\mathcal{W}_N\vx_0$, so the singular values of
$\mathcal{O}_N$ measure how visible each direction of the state space is, and the ratio
$\sigma_{\max}/\sigma_{\min}$ is a \emph{degree of unobservability}
\cite{kreneride2009}. A direction with $\sigma_j \approx 0$ is one whose error produces an
output signature below the sensor noise floor; no filter, however clever, will recover it from
that window.

\subsection{Nonlinear systems}
For \eqref{eq:prelim-proc}--\eqref{eq:prelim-meas} the corresponding notion is \emph{local weak
observability} \cite{hermannkrener1977}: the system is locally weakly observable at $\vx^\ast$ if
the differentials of the iterated output map
$\{h,\ h\circ f,\ h\circ f\circ f,\dots\}$ span $\real^{n_x}$ at $\vx^\ast$. Evaluating those
differentials along a trajectory reproduces \eqref{eq:prelim-obsmat} with the Jacobians
\eqref{eq:prelim-ekf-jac} frozen at the operating point, which is exactly the object an EKF
implicitly relies on. Two consequences matter here: observability is a \emph{local} and
\emph{input-dependent} property, so a battery model can be well observed during a hover segment
and nearly unobservable during a cruise plateau; and the linearised test is the right one to
apply, because the EKF's error dynamics are governed by the linearisation, not by the true map.

\subsection{The role of the OCV slope in battery estimation}
\label{sec:prelim-ocvslope}
For the cell model of \cref{ch:ecm}, $\vx = [\soc, v_1, v_2, h]^\T$ and the measurement Jacobian
is
\begin{equation}
\mH = \begin{bmatrix}
 \dfrac{\partial\ocv}{\partial\soc}(\soc,T) & -1 & -1 & M
\end{bmatrix},
\label{eq:prelim-battery-H}
\end{equation}
where $M$ is the hysteresis magnitude in volts. The SOC enters the measurement through one
scalar, the OCV slope $\partial\ocv/\partial\soc$, and through nothing else: the state is
otherwise a pure integrator, $\partial f_\soc/\partial\soc = 1$. Everything about SOC estimation
follows from this observation.
\begin{itemize}[nosep,leftmargin=*]
\item The slope is the \emph{observability gain}. A voltage error $\delta v$ maps to an SOC
  error $\delta\soc = \delta v / (\partial\ocv/\partial\soc)$. On the NMC811/graphite cell of
  this report the slope varies by a factor of six over the usable window and drops to
  $\SI{0.196}{\volt}$ near $\soc\approx0.90$; a \SI{4}{\milli\volt} model error there is worth
  two percentage points of SOC (\cref{sec:ecm-obs} quantifies this).
\item Where the slope vanishes, the SOC becomes unobservable from voltage and only Coulomb
  counting remains --- which is why chemistries with flat plateaux (LFP) are much harder than
  NMC, and why the benchmark scenarios deliberately start at $\soc_0 = 0.95$, on the flattest
  part of the curve.
\item A current-sensor bias is \emph{structurally} unobservable from a single voltage
  measurement unless it is added to the state: bias and SOC enter the output through the same
  channel, so the filter can only trade one against the other. Zhao, Duncan and
  Howey~\cite{zhao2017observability} prove this for the ECM with augmented sensor biases and show
  what extra excitation restores observability; Lin~\cite{lin2018sensorbias} derives the
  resulting steady-state SOC error as an explicit function of the bias and of the OCV slope.
  These two papers are the theoretical backdrop of the \code{current\_bias} scenario and of every
  disturbance-observer and unknown-input method in Part~VI.
\end{itemize}

\section{Consistency}
\label{sec:prelim-consistency}

A filter that reports a covariance makes a falsifiable claim, and the benchmark checks it. The
two standard statistics are due to Bar-Shalom \cite[\S 5.4]{barshalom2001}.

\begin{definition}[NIS and NEES]\label{def:prelim-nis}
The \emph{normalised innovation squared} and the \emph{normalised estimation error squared} at
sample $k$ are
\begin{equation}
\mathrm{NIS}_k = \veps_k^\T\mS_k\inv\veps_k, \qquad
\mathrm{NEES}_k = (\vx_k-\xplus_k)^\T(\Pplus_k)\inv(\vx_k-\xplus_k).
\label{eq:prelim-nisnees}
\end{equation}
\end{definition}

If the filter is consistent, $\mathrm{NIS}_k\sim\chi^2_{n_y}$ and
$\mathrm{NEES}_k\sim\chi^2_{n_x}$, so $\E{\mathrm{NIS}_k}=n_y$ and $\E{\mathrm{NEES}_k}=n_x$. Over
$N$ independent samples the time average $\bar{\mathrm{NIS}} = N\inv\sum_k\mathrm{NIS}_k$ has
$N\bar{\mathrm{NIS}}\sim\chi^2_{N n_y}$, so the two-sided acceptance region at confidence
$1-\alpha$ is
\begin{equation}
\bar{\mathrm{NIS}} \in
 \left[\ \frac{1}{N}\chi^2_{Nn_y}\!\left(\tfrac{\alpha}{2}\right),\ \
 \frac{1}{N}\chi^2_{Nn_y}\!\left(1-\tfrac{\alpha}{2}\right)\right],
\label{eq:prelim-nistest}
\end{equation}
which for the cell model ($n_y=1$) and a \SI{600}{\second} window at \SI{10}{\hertz}
($N=6000$) is $[0.964,\,1.036]$ at \SI{95}{\percent} --- a tight test, because a long record
leaves a badly tuned filter nowhere to hide.

NIS needs no ground truth and can therefore run on a real aircraft; NEES needs the true state and
is available only in simulation. A filter with $\bar{\mathrm{NIS}}\gg n_y$ is over-confident ---
the usual EKF failure --- and one with $\bar{\mathrm{NIS}}\ll n_y$ is conservative, wasting
information but safe. For a certified BMS the asymmetry is deliberate: a state-of-charge
indication with an honest or pessimistic interval is airworthy, an optimistic one is not.

Two caveats apply to everything in Part~XIII. First, \eqref{eq:prelim-nistest} assumes the
innovations are independent; when the model error is deterministic and load-correlated --- which
is exactly the situation on the electrochemical plant of \cref{ch:spm} --- the innovation
sequence is coloured, the effective sample size is far below $N$, and a failed test indicates
model error rather than covariance mis-tuning. Second, a filter can be consistent and useless
(report a huge covariance and a huge error) or inconsistent and accurate. Consistency is a
statement about the \emph{honesty} of $\mP$, not about the size of the error; the benchmark
therefore reports both, and neither alone.

\section{Cram\'er--Rao and posterior bounds}
\label{sec:prelim-crlb}

The benchmark reports achieved errors; it is worth knowing what the best achievable error is.
For an unbiased estimator of a deterministic parameter $\vtheta$ from data $\vy$, the
Cram\'er--Rao bound is $\operatorname{Cov}(\hat{\vtheta})\succeq \bm{J}\inv$ with Fisher
information $\bm{J} = \E{(\nabla_{\vtheta}\log p(\vy\mid\vtheta))(\nabla_{\vtheta}\log
p(\vy\mid\vtheta))^\T}$ \cite{vantrees1968}. For a random state the corresponding object is the
\emph{posterior} (Bayesian) Cram\'er--Rao lower bound, PCRLB, which for the model
\eqref{eq:prelim-proc}--\eqref{eq:prelim-meas} obeys the Riccati-like recursion of Tichavsk\'y,
Muravchik and Nehorai \cite{tichavsky1998pcrb},
\begin{equation}
\bm{J}_{k} = \left(\mQ_{k-1} + \mF_{k-1}\bm{J}_{k-1}\inv\mF_{k-1}^\T\right)\inv
  + \E{\mH_k^\T\mR_k\inv\mH_k},
\label{eq:prelim-pcrlb}
\end{equation}
where the expectation is over the true state trajectory and $\mF,\mH$ are the Jacobians
\eqref{eq:prelim-ekf-jac} evaluated at the truth. $\E{(\vx_k-\xhat_k)(\vx_k-\xhat_k)^\T}\succeq
\bm{J}_k\inv$ for any estimator.

Two remarks make \eqref{eq:prelim-pcrlb} useful rather than decorative. First, comparing it with
the EKF recursion \eqref{eq:prelim-kf-tu}--\eqref{eq:prelim-kf-mu} shows that they are the
\emph{same} recursion with the Jacobians evaluated at the truth rather than at the estimate: a
filter whose reported $\mP$ is below $\bm{J}\inv$ is provably lying. Second, for the cell model
the $(\soc,\soc)$ entry of \eqref{eq:prelim-pcrlb} is dominated by
$\E{(\partial\ocv/\partial\soc)^2}/\sigma_v^2$, so the achievable SOC accuracy scales as
$\sigma_v/(\partial\ocv/\partial\soc)$ divided by the square root of the number of independent
samples --- the same slope-driven statement as \cref{sec:prelim-ocvslope}, now with a constant.
At $\sigma_v=\SI{2}{\milli\volt}$, a slope of $\SI{0.8}{\volt}$ and \SI{10}{\hertz} sampling, a
one-minute window bounds the SOC standard deviation below about
$\SI{2}{\milli\volt}/(\SI{0.8}{\volt}\sqrt{600}) \approx 10^{-4}$, i.e.\ \num{0.01}\,\%. Measured
steady-state root-mean-square errors on the ECM plant are two orders of magnitude larger than
this, which is the quantitative statement that the benchmark is limited by \emph{model error},
not by sensor noise.

\section{Deterministic observers versus stochastic filters}
\label{sec:prelim-observers}

Part~VI of this report treats a family of estimators that never mention probability. It is worth
stating the relationship precisely, because the two viewpoints answer different questions and
the benchmark rewards them differently.

A Luenberger observer \cite{luenberger1964,luenberger1971} for the linear system
\eqref{eq:prelim-lin-proc}--\eqref{eq:prelim-lin-meas} is the copy of the plant driven by the
output error,
\begin{equation}
\xhat_k = \mF\xhat_{k-1} + \mG\vu_{k-1} + \mL\left(\vy_{k-1} - \mH\xhat_{k-1}\right),
\label{eq:prelim-luen}
\end{equation}
whose error $\tilde{\vx}_k = \vx_k-\xhat_k$ obeys the autonomous linear recursion
\begin{equation}
\tilde{\vx}_k = (\mF - \mL\mH)\,\tilde{\vx}_{k-1} + \vw_{k-1} - \mL\vv_{k-1}.
\label{eq:prelim-luen-err}
\end{equation}
The design problem is the placement of the eigenvalues of $\mF-\mL\mH$, which is solvable exactly
when $(\mF,\mH)$ is observable \cite{ackermann1972}. Structurally
\eqref{eq:prelim-luen} \emph{is} the Kalman filter with $\mL = \mF\mK$; the difference is
entirely in how $\mL$ is chosen and what is claimed about it:
\begin{itemize}[nosep,leftmargin=*]
\item The Kalman filter chooses $\mL$ to minimise the error covariance under an assumed noise
  model, and \emph{reports} that covariance. Its claim is statistical and is testable with
  \cref{def:prelim-nis}.
\item The observer chooses $\mL$ to shape the transient of \eqref{eq:prelim-luen-err} and
  typically claims a deterministic bound: exponential convergence with a given rate, or
  input-to-state stability with a given gain from $(\vw,\vv)$ to $\tilde{\vx}$
  \cite{khalil2002}. It reports no covariance, so it cannot be inconsistent --- and cannot be
  used to publish a confidence interval either.
\end{itemize}
The observer viewpoint dominates when the noise statistics are unknown or plainly non-Gaussian,
when the uncertainty is better described as a bounded disturbance than as a random variable
(sliding-mode and interval observers), or when a hard guarantee is required for certification.
The filter viewpoint dominates when the statistics are known and an uncertainty figure is part of
the deliverable. Both appear in this report on the same twelve scenarios, and the comparison is
one of its more interesting outcomes.

\section{Robustness}
\label{sec:prelim-robust}

``Robustness'' names three different problems in this report and they require three different
remedies.

\paragraph{Unknown noise statistics ($\mQ,\mR$ wrong).} The remedy is adaptation: estimate the
covariances online from the innovation sequence \cite{mehra1970,sagehusa1969,mohamed1999adaptive},
discount old data with a fading factor \cite{sorenson1971fading}, or treat the covariance as a
random variable with a conjugate prior and update it variationally
\cite{sarkka2009vb}. Part~IV is devoted to these.

\paragraph{Heavy-tailed or contaminated measurements.} A Gaussian likelihood assigns an
influence to a residual that grows linearly without bound, so a single \SI{50}{\milli\volt}
outlier moves a \SI{2}{\milli\volt}-noise filter by twenty-five standard deviations. The remedies
replace the quadratic loss with one whose influence function is bounded --- Huber's
$\rho$ \cite{huber1964,masreliez1975,karlgaard2007huber}, the correntropy criterion
\cite{chen2017mckf} --- or replace the Gaussian by a Student-$t$ whose tails accommodate the
outlier \cite{roth2013studentt}. Formally, the maximum-\emph{a-posteriori} update solves
\begin{equation}
\xplus_k = \argmin_{\vx}\ \rho\!\left(\vy_k - h(\vx,\vu_k);\,\mR_k\right)
  + \norm{\vx-\xminus_k}^2_{(\Pminus_k)\inv},
\label{eq:prelim-robustmap}
\end{equation}
with $\rho$ quadratic for the Kalman filter and sub-quadratic in the tails for the robust
variants; the iteratively reweighted least-squares solution of \eqref{eq:prelim-robustmap} is
literally a Kalman update with an inflated $\mR$.

\paragraph{Unknown model error.} Here the minimax view is the right one. The $H_\infty$
(minimax) filter makes no stochastic assumption at all and instead bounds the induced
$\ell_2$ gain from the disturbances to the estimation error,
\begin{equation}
\sup_{\vw,\vv,\vx_0\ne 0}\
\frac{\sum_k \norm{\vx_k-\xhat_k}^2_{\bm{S}}}
     {\norm{\vx_0-\xhat_0}^2_{\mP_0\inv} + \sum_k\left(\norm{\vw_k}^2_{\mQ\inv}+\norm{\vv_k}^2_{\mR\inv}\right)}
 < \gamma^2 ,
\label{eq:prelim-hinf}
\end{equation}
for a prescribed level $\gamma$ \cite{shen1997hinf,hassibi1996hinf,simon2006}. The resulting
recursion is a Kalman filter with an indefinite correction term in the Riccati equation; as
$\gamma\to\infty$ it reduces to the Kalman filter, and as $\gamma$ decreases it trades nominal
optimality for a guaranteed worst case. This is precisely the trade the SPM-plant benchmark of
\cref{ch:spm} probes, because there the model error is structured, deterministic and of the same
order as the sensor noise --- the regime in which the stochastic assumption is simply false.

\section{Symbols used throughout the report}
\label{sec:prelim-symbols}

\cref{tab:prelim-symbols} collects the recurring symbols with their units. Symbols local to a
single chapter are defined where they appear.

\begin{longtable}{@{}llp{78mm}@{}}
\caption{Principal symbols, with SI units where applicable.}\label{tab:prelim-symbols}\\
\toprule
Symbol & Unit & Meaning \\ \midrule
\endfirsthead
\toprule Symbol & Unit & Meaning \\ \midrule \endhead
\bottomrule \endfoot
\multicolumn{3}{@{}l}{\emph{Estimation}}\\
$k$ & --- & discrete sample index, $t_k = k\dt$ \\
$\dt$ & \si{\second} & sample period (\num{0.1} throughout the benchmark) \\
$\vx_k$ & --- & state vector, $\vx\in\real^{n_x}$ \\
$\vu_k$ & --- & input vector (true inputs and scheduling variables), $\vu\in\real^{n_u}$ \\
$\vy_k$ & --- & measurement vector, $\vy\in\real^{n_y}$ \\
$\vw_k,\vv_k$ & --- & process and measurement noise \\
$\mQ,\mR$ & --- & process- and measurement-noise covariances \\
$\xminus_k,\xplus_k$ & --- & prior and posterior state estimate \\
$\Pminus_k,\Pplus_k$ & --- & prior and posterior error covariance \\
$\mF,\mB,\mH$ & --- & Jacobians $\partial f/\partial\vx$, $\partial f/\partial\vu$, $\partial h/\partial\vx$ \\
$\mK_k$ & --- & Kalman gain; $\mL$ observer gain \\
$\veps_k$ & \si{\volt} & innovation $\vy_k-\yhat_k$ (voltage residual for the cell model) \\
$\mS_k$ & \si{\volt\squared} & innovation covariance \\
$\bm{Y},\bm\xi$ & --- & information matrix $\mP\inv$ and information vector $\mP\inv\xhat$ \\
$\bm{J}_k$ & --- & posterior Cram\'er--Rao information matrix \\
$\vtheta$ & --- & parameter vector for joint/dual estimation, $\vtheta=[Q,R_0]^\T$ \\
$\Pi,\mathcal{X}$ & --- & projection onto the feasible set and the feasible set itself \\[1mm]
\multicolumn{3}{@{}l}{\emph{Cell}}\\
$\soc$ & --- & state of charge, $\soc\in[0,1]$ \\
$v$ & \si{\volt} & terminal voltage (measurement) \\
$i$ & \si{\ampere} & cell current, discharge positive \\
$T$ & \si{\kelvin} & cell temperature ($T_{\mathrm{amb}}$ ambient, $T_{\mathrm{ref}}=\SI{298.15}{\kelvin}$) \\
$Q$ & \si{\ampere\hour} & present total cell capacity \\
$\ocv(\soc,T)$ & \si{\volt} & open-circuit voltage \\
$\partial\ocv/\partial\soc$ & \si{\volt} & OCV slope (observability gain) \\
$\partial U/\partial T$ & \si{\volt\per\kelvin} & entropic coefficient \\
$R_0$ & \si{\ohm} & series (ohmic) resistance \\
$R_1,R_2$ & \si{\ohm} & RC-branch resistances \\
$\tau_1,\tau_2$ & \si{\second} & RC-branch time constants \\
$v_1,v_2$ & \si{\volt} & RC-branch (diffusion/polarisation) voltages \\
$h$ & --- & one-state hysteresis variable, $h\in[-1,1]$ \\
$M$ & \si{\volt} & hysteresis magnitude; $\gamma$ hysteresis rate constant \\
$\eta_c$ & --- & coulombic efficiency on charge \\
$E_a$ & \si{\joule\per\mole} & Arrhenius activation energy \\
$C_{\mathrm{th}},R_{\mathrm{th}}$ & \si{\joule\per\kelvin}, \si{\kelvin\per\watt} & lumped heat capacity and thermal resistance \\
$\sigma_v,\sigma_i,\sigma_T$ & \si{\volt},\si{\ampere},\si{\kelvin} & sensor noise standard deviations \\[1mm]
\multicolumn{3}{@{}l}{\emph{Electrochemical model}}\\
$c_{\mathrm{n}},c_{\mathrm{p}}$ & \si{\mole\per\cubic\metre} & solid lithium concentration, negative/positive particle \\
$c_{\max}$ & \si{\mole\per\cubic\metre} & maximum solid concentration \\
$x,y$ & --- & electrode stoichiometries $c/c_{\max}$, negative/positive \\
$D_{\mathrm{n}},D_{\mathrm{p}}$ & \si{\square\metre\per\second} & solid diffusivities \\
$R_{\mathrm{n}},R_{\mathrm{p}}$ & \si{\metre} & particle radii \\
$a_{\mathrm{n}},a_{\mathrm{p}}$ & \si{\per\metre} & specific interfacial area $3\varepsilon/R$ \\
$J$ & \si{\mole\per\square\metre\per\second} & molar surface flux \\
$j_0$ & \si{\ampere\per\square\metre} & exchange-current density \\
$\eta_{\mathrm{n}},\eta_{\mathrm{p}}$ & \si{\volt} & kinetic (Butler--Volmer) overpotentials \\
$U_{\mathrm{n}},U_{\mathrm{p}}$ & \si{\volt} & electrode equilibrium potentials vs Li/Li$^+$ \\
$F$ & \SI{96485.33}{\coulomb\per\mole} & Faraday constant; $R_g=\SI{8.3145}{\joule\per\mole\per\kelvin}$ \\
\end{longtable}

\section{Summary}

The recursion \eqref{eq:prelim-ck}--\eqref{eq:prelim-bayes} is the problem; everything else in
this report is an approximation to it, chosen under a different combination of computational
budget, knowledge of the noise, and required guarantee. The Kalman filter solves it exactly in
the linear-Gaussian case and optimally among linear estimators otherwise; linearisation and
moment matching are the two ways to keep its form when the model is nonlinear; observability, and
in particular the OCV slope of \eqref{eq:prelim-battery-H}, decides what is recoverable at all;
consistency and the posterior Cram\'er--Rao bound decide whether a filter's own uncertainty claim
is honest and whether its error is close to the physical limit. \cref{ch:ecm} instantiates this
machinery on the equivalent-circuit cell model, and \cref{ch:spm} builds the electrochemical
plant against which the whole construction is tested.
